\chapter{Discusión general y conclusiones}
\label{chap:discusion}

\section{Síntesis de las dos aristas}

Los dos capítulos centrales de esta tesis abordan la misma pregunta —¿por qué la coincidencia de fragilidad bancaria y riesgo de cola del crecimiento eleva el spread soberano más de lo que la suma de ambos factores explicaría?— desde ángulos que se validan con estándares distintos, y llegan a una respuesta convergente.

El Capítulo~\ref{chap:paper2} muestra, con un panel de cinco economías latinoamericanas y controles macroeconómicos completos, que el coeficiente de interacción entre $JLoss$ y $GaR$ sobre el EMBI es negativo y significativo bajo inferencia estándar ($\hat\theta=-0{,}338$, $t=-2{,}22$), estable frente a una batería exigente de robustez, y corroborado por un modelo de umbral que no impone la forma funcional de la interacción. El mismo capítulo documenta, con la misma honestidad, que esta regularidad se atenúa y pierde solidez fuera del núcleo latinoamericano, y que su significancia puntual —aunque no su signo— es sensible al reducido número de \textit{clusters} de país disponible bajo inferencia robusta.

El Capítulo~\ref{chap:paper1} muestra que esta complementariedad no es una regularidad estadística sin anclaje estructural: se deriva formalmente de un modelo de competencia bancaria en el que la estructura de mercado gobierna, en cadena, la fragilidad individual, la fragilidad de cola conjunta del sistema, el canal de \textit{credit crunch} hacia el crecimiento, y finalmente el spread soberano. El modelo genera, además, una predicción que el capítulo empírico por sí solo no habría podido anticipar sin él: que la concentración bancaria amplifica la complementariedad. Esa predicción distintiva se confirma con el signo correcto en cuatro de cinco proxies de concentración y competencia sobre un panel de once economías emergentes, con la salvedad honesta —declarada explícitamente en ambos capítulos— de que el poder estadístico disponible hoy es moderado.

La convergencia de ambas aristas es el resultado más defendible de esta tesis: no un coeficiente puntual preciso en ninguno de los dos capítulos, sino una cadena \textbf{teoría–medición–evidencia} que se sostiene en cada uno de sus eslabones. Un modelo estructural predice la complementariedad y su amplificación por concentración; un panel de datos reales, construido de forma independiente y homogénea entre países, confirma el signo de ambas predicciones; y las limitaciones de poder estadístico de la evidencia —lejos de ocultarse— quedan documentadas como la frontera actual, no como el techo definitivo, de lo que este programa de investigación puede afirmar.

\section{Implicancias de política conjuntas}

La lectura de política que emerge de ambos capítulos es más precisa que la que cualquiera de los dos habría entregado por separado. El Capítulo~\ref{chap:paper2} ya establece que la regulación orientada a la resiliencia bancaria —capital, liquidez, provisiones contracíclicas— reduce la exposición del soberano a los escenarios macroeconómicos adversos, y que la fragilidad bancaria y el riesgo de cola del crecimiento deben monitorearse conjuntamente. El Capítulo~\ref{chap:paper1} añade una dimensión que la sola evidencia empírica del spread soberano no puede entregar: la estructura de mercado bancario misma es un objeto de política. El modelo muestra que el mercado que minimiza el riesgo de cola \textit{sistémico} está más competido que el que minimiza el riesgo de quiebra \textit{individual} de un banco representativo (Proposición~\ref{prop:shift}) —una discrepancia que, en el marco de la organización industrial bancaria estándar, permanece invisible—, y que la concentración amplifica la transmisión de cualquier fragilidad remanente al riesgo soberano. Una política de competencia bancaria calibrada únicamente sobre indicadores de salud individual de los bancos, en consecuencia, puede dejar al sistema en una configuración subóptimamente concentrada desde la perspectiva del riesgo sistémico y de su costo fiscal contingente. La regulación macroprudencial y la política de competencia bancaria no son, a la luz de esta tesis, instrumentos independientes: ambas gobiernan el mismo parámetro profundo.

\section{Limitaciones comunes}

Ambos capítulos comparten dos limitaciones de fondo que conviene reconocer de manera unificada. Primera, la cobertura de países con balances bancarios homogéneos, $JLoss$ construido y EMBI real sigue siendo, en el estado actual del \textit{pipeline}, el cuello de botella que limita tanto la generalización de la evidencia empírica del Capítulo~\ref{chap:paper2} fuera de América Latina como el poder estadístico de la prueba de amplificación por concentración del Capítulo~\ref{chap:paper1}. Segunda, ninguno de los dos capítulos cierra completamente el problema de simultaneidad del \textit{doom loop}: el Capítulo~\ref{chap:paper2} lo aborda con una batería de estrategias de identificación causal que confirman la dirección banco→soberano del canal de nivel pero no una magnitud causal precisa de la interacción; el Capítulo~\ref{chap:paper1} lo aborda con un bloque soberano de equilibrio parcial, cuyo cierre dinámico de dos períodos queda como agenda pendiente.

\section{Agenda de investigación futura}

De ambos capítulos se desprende una agenda común y priorizada. En primer lugar, la ampliación del panel de economías con $JLoss$, $GaR$ y EMBI real construidos de forma homogénea es la vía directa para resolver simultáneamente la generalización del Capítulo~\ref{chap:paper2} y la precisión estadística del Capítulo~\ref{chap:paper1}: varios extractores de balances bancarios ya están escritos y solo pendientes de ejecución, y la extensión de la cobertura de EMBI real más allá de 2014 para las economías no latinoamericanas es un problema de acceso a datos, no de diseño. En segundo lugar, el tratamiento formal de la dinámica del \textit{doom loop} —mediante estimadores GMM dinámicos en el capítulo empírico, y mediante una extensión de dos períodos del bloque soberano en el capítulo teórico— permitiría cerrar la brecha entre la identificación de nivel, hoy razonablemente establecida, y la identificación causal de la interacción, todavía abierta. En tercer lugar, el microfundamento completo de la correlación de activos $\rho(n)$ más allá del caso de competencia espacial de Salop —hoy postulada en la versión base del modelo— fortalecería la robustez teórica de la Proposición~\ref{prop:shift}. Estas tres líneas no son independientes: cada una, al completarse, refuerza directamente la evidencia que la otra arista de esta tesis ya ha comenzado a construir.

\section{Conclusión}

Esta tesis sostiene, con la evidencia y el rigor disponibles hoy, que la fragilidad bancaria sistémica y el riesgo a la baja del crecimiento no son determinantes independientes del riesgo soberano en economías emergentes, sino complementos cuya coincidencia amplifica el spread soberano por encima de la suma de sus efectos individuales —y que esa amplificación misma crece con la concentración del sistema bancario. La arista empírica establece esta complementariedad como una regularidad robusta dentro de su ámbito de validez declarado; la arista teórica la deriva desde primeros principios de un modelo de competencia bancaria y extiende la predicción a un terreno —la amplificación por concentración— que la literatura de organización industrial bancaria estándar no anticipa. Juntas, delimitan con precisión tanto lo que puede afirmarse hoy como lo que constituye, honestamente, la agenda pendiente.
